\section*{Part II: Selecting the Right Geography}
\addcontentsline{toc}{section}{Part II: Selecting the Right Geography}


\subsection{7. A city's talent ecology often matters more than its land price}

Faculty and researchers make life decisions, not merely commuting decisions. A location near research clusters, technically sophisticated employers, schools for children, healthcare, airports and professional communities can improve recruiting and collaboration. The price paid for this access can be rational. Conversely, a low-price parcel far from universities and industry may require higher salaries, daily transport, housing subsidies and travel budgets for decades.

Model \textbf{total faculty and student friction}, not simply kilometres to a city. A cheap site can create a permanent operational tax.


\subsection{8. Begin with a defined research and academic niche}

A materials lab, AI research center, biotechnology institute, design school and residential university require different ecosystems. For a STEM/AI-1st institution, relevant location attributes may include software/AI employers, quantitative faculty, student supply, compute and electrical infrastructure, academic partners, and international accessibility. For wet labs, reliable water, effluent treatment, specialized contractors and permitting may dominate. Location weights should be driven by activities-not borrowed from a generic residential property article.


\subsection{9. City-level longlist: a process, not a verdict}

For a Maharashtra-based founder, candidate regions could include metropolitan Mumbai, Pune's wider research and industrial corridor, Nashik, Nagpur and selected emerging locations. These are candidates for research, not recommended land purchases. Compare them with at least 1 outside-state or foreign reference market if an internationally oriented campus is genuinely in scope. The result may show that a leased urban research center and a later greenfield campus are different locations-and need not be forced together.

Do not assume that either India's cheapest land or the globally most prestigious city maximizes institutional value.


\subsection{10. A location feasibility hierarchy}

Apply \textbf{hard gates} before any weighted ranking: can the proposed institutional use legally occur; can the required land tenure be obtained; is physical access lawful; is water and power viable; is environmental or flood exposure manageable; can a financeable minimum phase be built; and is there a plausible route for faculty and students to reach it? A parcel failing a hard gate should not be rescued by a high aggregate score.

Then evaluate talent access, customers/partners, connectivity, livability, total development cost, sponsor/regulatory context, climate resilience and expansion potential. Use weighted scores only for feasible options.


\subsection{11. A worked geographic scoring exercise}

\textbf{[Model]} Compare 3 \textbf{hypothetical} locations A, B, and C. Score each from 1 (poor) to 5 (excellent), after hard legal gates. Example weights are strategic judgments rather than market data.

\begin{table}[t]
\centering
\small
\caption{Worked Geographic Scoring Matrix (Illustrative Model)}
\label{tab:ch34_geo_scoring}
\begin{tabularx}{\columnwidth}{X r r r r}
\toprule
\textbf{Dimension} & \textbf{Weight} & \textbf{A} & \textbf{B} & \textbf{C} \\
\midrule
Faculty / research talent & 25\% & 5 & 4 & 3 \\
Industry / student access & 20\% & 5 & 4 & 3 \\
All-in long-run site cost & 20\% & 2 & 3 & 5 \\
Mobility and livability & 15\% & 4 & 4 & 3 \\
Infrastructure / resilience & 10\% & 4 & 4 & 3 \\
Expansion optionality & 10\% & 2 & 4 & 5 \\
\midrule
\textbf{Weighted score / 5} & \textbf{100\%} & \textbf{3.85} & \textbf{3.80} & \textbf{3.60} \\
\bottomrule
\end{tabularx}
\end{table}
A narrowly leads B on these assumed weights, but the difference is small despite radically different economics. That is useful: it shows the decision depends on mission priorities and hard-to-measure factors. Apply sensitivity tests, cost estimates, and actual evidence. The score does not license acquisition.


\subsection{12. Sensitivity analysis is mandatory}

Change the weights materially: what if faculty access rises from 25% to 35%, or if cost becomes 30%? Ask whether the rank changes. If it does, the winning site is not robust. Compare the cheapest feasible micro-markets within each winning region rather than averaging over a whole city.

Also test a \textbf{2-site strategy}-research nucleus near a talent cluster and a later residential campus in a nearby lower-cost locality. Such separation introduces travel and coordination costs, but may preserve quality at lower capital outlay.


\subsection{13. Travel-time is often the real distance}

For daily students and commuters, actual travel times, variability, last-mile access, safe public transport, parking and off-peak service matter more than road kilometres. For international faculty, airport access plus predictable onward travel matters. For residential campuses, connections to hospitals, shopping, schools and family employment matter.

Conduct timed visits at commuting peaks and during adverse weather, and interview prospective residents and staff. Model the \textbf{effective annual cost of time}. Even an unpriced hour of friction repeated across hundreds of faculty and staff becomes a major institutional burden.


\subsection{14. Understanding the land market}

Ask brokers for actual closed-sale evidence, not only asking prices; compare per legally usable square metre or per permissible built-up area, not per quoted acre. Check parcel shape, road width, permitted floor area, setbacks, slope, soil, flood zone, land assembly requirements, and utility connections. A narrow plot with a beautiful per-acre price may produce a poor building footprint. A parcel with an access dispute may have nearly zero educational use until resolved.

Use multiple independent brokers and professional valuation where material. Record who benefits from the transaction, including intermediaries and related parties.


\subsection{15. The right site may require no purchase}

Before commissioning acquisition counsel, test a short residency or lab pilot in the target region. Measure faculty willingness to travel, local student demand, actual utility reliability, vendor availability, and operating cost. A \textnormal{₹}2-lakh pilot that invalidates a \textnormal{₹}20-crore site thesis has extremely high decision value.

Small reversible tests can reveal more than months of speculative map analysis, provided the tests mirror the actual intended use.
